% Part: proof-theory
% Chapter: cut-elimination
% Section: ce-topmost

\documentclass[../../../include/open-logic-section]{subfiles}

\begin{document}

\olfileid{pt}{cut}{top}

\olsection{Removing Topmost Cuts}

\begin{defn}
Consider !!a{proof}~$\pi$ ending in a \CutCS{} rule,
\[
\AxiomC{}
\RightLabel{$\pi_1$}
\Deduce$\Gamma \fCenter \Delta, !A$
\AxiomC{}
\RightLabel{$\pi_2$}
\Deduce$!A, \Gamma \fCenter \Delta$
\RightLabel{\CutCS}
\BinaryInf$\Gamma \fCenter \Delta$
\DisplayProof
\]
and containing no other \CutCS{} inferences. The \emph{cut height}
of~$\pi$ is $\cheight{\pi} = \pheight{\pi_1} + \pheight{\pi_2}$. The
\emph{cut rank} of~$\pi$ is $\cutrank{\pi} = \depth{!A}$.
\end{defn}

\begin{lem}\ollabel{lem:cut-adm-G3c}
The \CutCS{} rule is admissible in~\Log{G3c}. In other words, if $\pi$
is !!a{proof} ending in a single~\CutCS{} and otherwise using only rules
of~$\Log{G3c}$, then there is !!a{proof}~$\pi'$ of the same end-sequent in~\Log{G3c}.
\end{lem}

\begin{proof}
The proof is by double induction on cut !!{height} and rank. The
!!{proof}~$\pi$ ends in:
\[
\AxiomC{}
\RightLabel{$\pi_1$}
\Deduce$\Gamma \fCenter \Delta, !A$
\AxiomC{}
\RightLabel{$\pi_2$}
\Deduce$!A, \Gamma \fCenter \Delta$
\RightLabel{\CutCS}
\BinaryInf$\Gamma \fCenter \Delta$
\DisplayProof
\]

Induction basis: $\cheight{\pi} = 0$ and $\cutrank{\pi} = 0$. Since
$\cheight{\pi} = \pheight{\pi_1} + \pheight{\pi_2} = 0$, both $\Gamma
\Sequent \Delta, !A$ and $!A, \Gamma \Sequent \Delta$ are axioms.
There are two cases: (a)~either $!A$ is not principal in one of them,
or (b)~$!A$ is principal in both of them. In both cases we'll show
below that $\Gamma \Sequent \Delta$ must be an axiom (and we don't
need the inductive hypothesis to do that).

We'll now distinguish a number of possible cases according to whether
$\pi_1$ and $\pi_2$ are axioms or end in an inference, and whether
$!A$~is principal in that inference or not. Cases A and~B cover the
induction basis. In the inductive step we assume $\cheight{\pi} > 0$
or $\cutrank{\pi} > 0$ (or both). We can use the inductive hypothesis,
namely that the last~\CutCS{} can be removed from any !!{proof} which
ends in a single~\CutCS{} which either has cut rank $= \cutrank{\pi}$
and a cut !!{height} $< \cheight{\pi}$, or has cut rank~$<
\cutrank{\pi}$. The case where $\cheight{\pi} = 0$ (both premises are
axioms) are covered by cases~A and~B. The cases where $\cheight{\pi} >
0$ are covered by cases C--D.

\paragraph{Case A:}
Some $!B$ occurs in both $\Gamma$ and~$\Delta$. 
Then $\Gamma \Sequent \Delta$ is an axiom (if $!B$~is atomic), or has
!!a{proof}~$\pi'$ in \Log{G3c} without \CutCS{} by
\olref[seq][ex]{prop:G3c-axioms}. This covers case~(a) of the
induction basis, since if $!A$ is not principal in $\Gamma
\Sequent \Delta, !A$ or not principal in $!A, \Gamma \Sequent \Delta$
and these are axioms, then some other atomic~$!B$ must occur in both
$\Gamma$ and~$\Delta$. It also covers the inductive step if one of the
premises is an axiom in which $!A$ is not principal (even if the other
premise is the conclusion of a rule).

\paragraph{Case B:} Both premises are axioms, and $!A$ is principal
and hence atomic. Consider the left premise, $\Gamma \Sequent
\Delta, !A$. Since $!A$ is principal, it must occur in $\Gamma$
(otherwise $\Gamma \Sequent \Delta, !A$ wouldn't be an axiom). Similarly,
$!A$ must occur in $\Delta$, or else $!A, \Gamma \Sequent \Delta$
wouldn't be an axiom. Thus $!A$ is atomic and occurs in both $\Gamma$
and $\Delta$, i.e., $\Gamma \Sequent \Delta$ is an axiom. This covers
case~(b) of the induction basis.

\begin{center}
Alternative cases A and B
\end{center}

\paragraph{Case A:} One of the premises of the cut is an axiom of the
form $!C, \Gamma' \Sequent \Delta', !C$ with $!C$~atomic. If $!C$ is
not the cut !!{formula}~$!A$, then $!C$ must occur in both $\Gamma$
and~$\Delta$. In that case, $\Gamma \Sequent \Delta$ is an axiom and
we can take it as~$\pi'$. If $!C \ident !A$ is the cut !!{formula}
then either (if the left premise is an axiom) $!A \in \Gamma$ and
$\Gamma = \Gamma', !A$, or (if the right premise is an axiom) $!A \in
\Delta$ and $\Delta = \Delta', !A$. In the first case, $\pi_2$ is
!!a{proof} of $!A, !A, \Gamma' \Sequent \Delta$ and by admissibility
of \LeftR{\Contraction} (\olref[seq][inv]{prop:G3c-cont-adm}) there is
!!a{proof} of $\Gamma \Sequent \Delta$. In the second case, $\pi_2$ is
!!a{proof} of $\Gamma \Sequent \Delta', !A, !A$ and by admissibility
of \RightR{\Contraction} we get !!a{proof} of $\Gamma \Sequent
\Delta$.

\paragraph{Case B:} One of the premises is an axiom of the
form~$\lfalse, \Gamma' \Sequent \Delta'$. If the left premise is such
an axiom, then $\lfalse \in \Gamma$ and $\Gamma \Sequent \Delta$ is
also an axiom. If the right premise is such an axiom, then either
$\lfalse \in \Gamma$ (in which case $\Gamma \Sequent \Delta$ again is
an axiom), or $!A \ident \lfalse$. In the latter case, $\pi_1$ is
!!a{proof} of $\Gamma \Sequent \Delta, \lfalse$. We can turn it into
!!a{proof} of~$\Gamma \Sequent \Delta$ by deleting one copy
of~$\lfalse$ from the consequent of every sequent in~$\pi_1$.
(Exercise: verify this by induction on~$\pheight{\pi_1}$.)

Now suppose that neither case~A nor case~B applies. In the remaining
cases, $\cheight{\pi} > 0$, i.e., at least one of the premises is the
conclusion of an inference.

\paragraph{Cases C and D: Permuting cuts.}
The possibilities considered in cases~C and~D all have a common
pattern. We begin with !!a{proof} ending in~\CutCS{} where one premise
is obtained by a rule~$R$ in which the cut !!{formula}~$!A$ is not
principal (some other !!{formula}~$!D$ is). In this !!{proof},
\CutCS{} follows the rule~$R$, and schematically looks as follows:
\[
\AxiomC{}
\RightLabel{$\pi_1$}
\Deduce$\Gamma \fCenter \Delta', !D, !A$
\AxiomC{}
\RightLabel{$\pi_3$}
\Deduce$!A, \Gamma, \Pi \fCenter \Delta', \Lambda$
\RightLabel{$R$}
\UnaryInf$!A, \Gamma \fCenter \Delta', !D$
\RightSubproofLabel{$\pi_2$}
\RightLabel{\CutCS}
\BinaryInf$\Gamma \fCenter \Delta', !D$
\DisplayProof
\]
This only covers the case where $R$~is a right rule with one premise,
$A$~is not the principal !!{formula} of~$R$, and the !!{formula}~$!D$ which
\emph{is} principal occurs in the succedent of the right premise
of~\CutCS. If $!D$~occurs in the antecedent (i.e., $R$ is a left rule)
or in the left premise of~\CutCS, it would look different, of course.
The !!{formula}s in $\Pi$ and~$\Lambda$ represent the active
!!{formula}s of rule~$R$.

We turn this order around and consider !!a{proof}
where $R$ follows~\CutCS{}:
\[
\AxiomC{}
\RightLabel{$\pi_1'$}
\Deduce$\Gamma, \Pi \fCenter \Delta', !D, \Lambda, !A$
\AxiomC{}
\RightLabel{$\pi_3'$}
\Deduce$!A, \Gamma, \Pi \fCenter \Delta', !D, \Lambda$
\RightLabel{\CutCS}
\BinaryInf$\Gamma, \Pi \fCenter \Delta', !D, \Lambda$
\RightLabel{$R$}
\UnaryInf$\Gamma, \fCenter \Delta', !D, !D$
\DisplayProof
\]
We use !!{height}-preserving weakening to obtain $\pi_1'$ and $\pi_3'$
from $\pi_1$ and $\pi_3$, respectively.

Because we switch the order of \CutCS{} and~$R$, this is referred to
as ``permuting a cut upwards across~$R$.'' Doing so decreases the
height of the !!{proof} ending in the right premise of~\CutCS, and so
decreases the cut !!{height}. That is, the sub-!!{proof}s ending in the
new \CutCS{} can be assumed to have !!{height} no larger than
$\pi_1$ and $\pi_3$, so the cut !!{height} is reduced (we've removed
one inference on the right). Now the induction hypothesis applies,
which tells us that the \CutCS{} inference can be removed. Finally, we
remove the extra copy of~$!D$ using the admissibility
of~$\RightR{\Contraction}$.

\paragraph{Case C:}
At least one of $\pi_1$ and $\pi_2$ ends in a one-premise rule~$R$ in
which $!A$~is not principal. There are a total of 18 cases, depending
on whether $!A$~is not principal in the left or the right premise
of~\CutCS, and depending on which of the nine one-premise rules
(\LeftR{\lnot}, \RightR{\lnot}, \RightR{\lor}, \LeftR{\land},
\RightR{\lif}, and the four quantifier rules) the rule~$R$ is. We only
consider two examples: where $!A$~is not principal in the last
inference of~$\pi_2$, which ends in~$\RightR{\lif}$ or
in~\LeftR{\lexists}. The other cases are analogous and left as
exercises.

Suppose $\pi$ ends in:
\[
\AxiomC{}
\RightLabel{$\pi_1$}
\Deduce$\Gamma \fCenter \Delta', !B \lif !C, !A$
\AxiomC{}
\RightLabel{$\pi_3$}
\Deduce$!A, !B, \Gamma \fCenter \Delta', !C$
\RightLabel{$\RightR{\lif}$}
\UnaryInf$!A, \Gamma \fCenter \Delta', !B \lif !C$
\RightSubproofLabel{$\pi_2$}
\RightLabel{\CutCS}
\BinaryInf$\Gamma \fCenter \Delta', !B \lif !C$
\DisplayProof
\]
where $\Delta = \Delta', !B \lif !C$. We apply !!{height}-preserving
weakening (\olref[seq][adm]{prop:weak-G3c-adm}) to $\pi_1$ to obtain
!!a{proof} $\pi_1'$ of $!B, \Gamma \fCenter \Delta', !B \lif !C, !C,
!A$ (we weaken with $!B$ on the left and $!C$ on the right).
Similarly, we apply \olref[seq][adm]{prop:weak-G3c-adm}
to~$\pi_3$ to get !!a{proof}~$\pi_3'$ of $!A, !B, \Gamma \fCenter
\Delta', !B \lif !C, !C$ (by weakening with $!B \lif !C$ on the right).
The induction hypothesis applies to
\[
\AxiomC{}
\RightLabel{$\pi_1'$}
\Deduce$!B, \Gamma \fCenter \Delta', !B \lif !C, !C, !A$
\AxiomC{}
\RightLabel{$\pi_3'$}
\Deduce$!A, !B, \Gamma \fCenter \Delta', !B \lif !C, !C$
\RightLabel{\CutCS}
\BinaryInf$!B, \Gamma \fCenter \Delta', !B \lif !C, !C$
\DisplayProof
\]
with cut !!{formula}~$!A$: This !!{proof} has the same cut rank and a
lower cut !!{height} than~$\pi$. This gives us !!a{proof}~$\pi_4$ in
\Log{G3c} without \CutCS{} of~$!B, \Gamma \fCenter \Delta', !B \lif
!C, !C$. Apply $\RightR{\lif}$ to get !!a{proof} of $\Gamma \fCenter
\Delta', !B \lif !C, !B \lif !C$:
\[
\AxiomC{}
\RightLabel{$\pi_4$}
\Deduce$!B, \Gamma \fCenter \Delta', !B \lif !C, !C$
\RightLabel{$\RightR{\lif}$}
\UnaryInf$\Gamma \fCenter \Delta', !B \lif !C, !B \lif !C$
\DisplayProof
\]
To get !!a{proof}~$\pi'$ of $\Gamma \Sequent \Delta$, i.e., $\Gamma
\fCenter \Delta', !B \lif !C$, apply
\olref[seq][inv]{prop:G3c-cont-adm} (the admissibility of
contraction).

Now suppose $\pi$ ends in:
\[
\AxiomC{}
\RightLabel{$\pi_1$}
\Deduce$\lexists[x][!B(x)], \Gamma' \fCenter \Delta, !A$
\AxiomC{}
\RightLabel{$\pi_3$}
\Deduce$!A, !B(c), \Gamma' \fCenter \Delta$
\RightLabel{$\LeftR{\lexists}$}
\UnaryInf$!A, \lexists[x][!B(x)], \Gamma' \fCenter \Delta$
\RightSubproofLabel{$\pi_2$}
\RightLabel{\CutCS}
\BinaryInf$\Gamma' \fCenter \Delta$
\DisplayProof
\]
We again apply the inductive hypothesis to
\[
\AxiomC{}
\RightLabel{$\pi_1[!B(c) \Sequent {}]$}
\Deduce$\lexists[x][!B(x)], !B(c), \Gamma' \fCenter \Delta, !A$
\AxiomC{}
\LeftLabel{$\pi_3[\lexists[x][!B(x)] \Sequent {}]$}
\Deduce$!A, \lexists[x][!B(x)], !B(c), \Gamma' \fCenter \Delta$
\RightLabel{\CutCS}
\BinaryInf$\lexists[x][!B(x)], !B(c), \Gamma' \fCenter \Delta$
\RightLabel{\LeftR{\lexists}}
\UnaryInf$\lexists[x][!B(x)], \lexists[x][!B(x)], \Gamma' \fCenter \Delta$
\DisplayProof
\]
Note that the eigenvariable condition is satisfied, since the
eigenvariable condition on~\LeftR{\lexists} in~$\pi_2$ guarantees that
$c$ does not occur in $\lexists[x][!B(x)]$, $\Gamma'$, or~$\Delta$.
Finally, apply \olref[seq][inv]{prop:G3c-cont-adm}.


\paragraph{Case D:}
At least one of $\pi_1$ and $\pi_2$ ends in a two-premise rule~$R$ in
which $!A$ is not principal. There are 6 cases. We consider the case
where $!A$ is not principal in the last inference of~$\pi_2$, which
ends in~$\RightR{\land}$; the other cases are analogous.

Assume $\pi$ ends in:
\[
\AxiomC{}
\RightLabel{$\pi_1$}
\Deduce$\Gamma \fCenter \Delta', !B \land !C, !A$
\AxiomC{}
\RightLabel{$\pi_3$}
\Deduce$!A, \Gamma \fCenter \Delta', !B$
\AxiomC{}
\RightLabel{$\pi_4$}
\Deduce$!A, \Gamma \fCenter \Delta', !C$
\RightLabel{$\RightR{\land}$}
\BinaryInf$!A, \Gamma \fCenter \Delta', !B \land !C$
\RightSubproofLabel{$\pi_2$}
\RightLabel{\CutCS}
\BinaryInf$\Gamma \fCenter \Delta$
\DisplayProof
\]
The inductive hypothesis applies to the !!{proof}s
\[
\AxiomC{}
\RightLabel{$\pi_1'$}
\Deduce$\Gamma \fCenter \Delta', !B \land !C, !B, !A$
\AxiomC{}
\RightLabel{$\pi_3'$}
\Deduce$!A, \Gamma \fCenter \Delta', !B \land !C, !B$
\RightLabel{\CutCS}
\BinaryInf$\Gamma \fCenter \Delta', !B \land !C, !B$
\DisplayProof
\]
and
\[
\AxiomC{}
\RightLabel{$\pi_1''$}
\Deduce$\Gamma \fCenter \Delta, !B \land !C, !C, !A$
\AxiomC{}
\RightLabel{$\pi_4'$}
\Deduce$!A, \Gamma \fCenter \Delta', !B \land !C, !C$
\RightLabel{\CutCS}
\BinaryInf$\Gamma \fCenter \Delta', !B \land !C, !C$
\DisplayProof
\]
Here, the \Log{G3c} !!{proof}s $\pi_1'$ and $\pi_1''$ are obtained by
!!{height}-preserving weakening (\olref[seq][adm]{prop:weak-G3c-adm})
from~$\pi_1$ (once weakening on the right with $!B$, and once with
$!C$). The !!{proof}s $\pi_3'$ and~$\pi_4'$ are obtained from $\pi_3$
and~$\pi_4$ also by !!{height}-preserving weakening from $\pi_3$
and~$\pi_4$, respectively (by weakening with $!B \land !C$ on the
right).

The inductive hypothesis gives us !!{proof}s~$\pi_5$ and~$\pi_6$ in
\Log{G3c} of $\Gamma \fCenter \Delta', !B \land !C, !B$ and $\Gamma
\fCenter \Delta', !B \land !C, !C$. We combine these into !!a{proof}
of $\Gamma \Sequent \Delta', !B \land !C, !B \land !C$ using the
rule~$\RightR{\land}$:
\[
\AxiomC{}
\RightLabel{$\pi_5$}
\Deduce$\Gamma \fCenter \Delta', !B \land !C, !B$
\AxiomC{}
\RightLabel{$\pi_6$}
\Deduce$\Gamma \fCenter \Delta', !B \land !C, !C$
\RightLabel{\RightR{\land}}
\BinaryInf$\Gamma \fCenter \Delta, !B \land !C, !B \land !C$
\DisplayProof
\]
To get !!a{proof} of $\Gamma \Sequent \Delta$, i.e., $\Gamma \fCenter
\Delta', !B \land !C$, apply \olref[seq][inv]{prop:G3c-cont-adm} (the
admissibility of contraction).

\paragraph{Case E: Reducing cuts.}
The final case concerns the situation where both $\pi_1$ and $\pi_2$
end in inferences in which $!A$ is principal. There are six cases,
one for each possible !!{main operator} of~$!A$.

In each case, we turn a \CutCS{} with cut !!{formula}~$!A$ into one or more
cuts on sub-!!{formula}s of~$!A$, namely on the active !!{formula}s of
the inferences of which~$!A$ is the principal !!{formula}. These cases
are therefore often referred to as ``reduction'' or ``conversion'' of
the cut inference.

\begin{enumerate}
\item If $!A \ident \lnot !B$ then $\pi$ ends in
\[
\AxiomC{}
\RightLabel{$\pi_1'$}
\Deduce$!B, \Gamma \fCenter \Delta$
\RightLabel{\RightR{\lnot}}
\UnaryInf$\Gamma \fCenter \Delta, \lnot !B$
\LeftSubproofLabel{$\pi_1$}
\AxiomC{}
\RightLabel{$\pi_2'$}
\Deduce$\Gamma \fCenter \Delta, !B$
\RightLabel{\LeftR{\lnot}}
\UnaryInf$\lnot !B, \Gamma \fCenter \Delta$
\RightSubproofLabel{$\pi_2$}
\RightLabel{\CutCS}
\BinaryInf$\Gamma \fCenter \Delta$
\DisplayProof
\]
By inductive hypothesis, \CutCS{} can be eliminated from
\[
\AxiomC{}
\RightLabel{$\pi_2'$}
\Deduce$\Gamma \fCenter \Delta, !B$
\AxiomC{}
\RightLabel{$\pi_1'$}
\Deduce$!B, \Gamma \fCenter \Delta$
\RightLabel{\CutCS}
\BinaryInf$\Gamma \fCenter \Delta$
\DisplayProof
\]

\item
If $!A \ident (!B \lor !C)$ then $\pi$ ends in
\[
\AxiomC{}
\RightLabel{$\pi_1'$}
\Deduce$\Gamma \fCenter \Delta, !B, !C$
\RightLabel{\RightR{\lor}}
\UnaryInf$\Gamma \fCenter \Delta, !B \lor !C$
\LeftSubproofLabel{$\pi_1$}
\AxiomC{}
\RightLabel{$\pi_2'$}
\Deduce$!B, \Gamma \fCenter \Delta$
\AxiomC{}
\RightLabel{$\pi_2''$}
\Deduce$!C, \Gamma \fCenter \Delta$
\RightLabel{\LeftR{\lor}}
\BinaryInf$!B \lor !C, \Gamma \fCenter \Delta$
\RightSubproofLabel{$\pi_2$}
\RightLabel{\CutCS}
\BinaryInf$\Gamma \fCenter \Delta$
\DisplayProof
\]
Consider the !!{proof} ending in \CutCS,
\[
\AxiomC{}
\RightLabel{$\pi_1'$}
\Deduce$\Gamma \fCenter \Delta, !B, !C$
\AxiomC{}
\RightLabel{$\pi_2'''$}
\Deduce$!C, \Gamma \fCenter \Delta, !B$
\RightLabel{\CutCS}
\BinaryInf$\Gamma \fCenter \Delta, !B$
\DisplayProof
\]
where $\pi_2'''$ is obtained from $\pi_2''$ by !!{height}-preserving weakening.
It has cut rank $< \cutr{\pi}$ as $\depth{!C} < \depth{!B \lor !C}$, so the
inductive hypothesis yields !!a{proof} in \Log{G3c} of $\pi_1''$ of
$\Gamma \fCenter \Delta, !B$. The !!{proof} ending in \CutCS,
\[
\AxiomC{}
\RightLabel{$\pi_1''$}
\Deduce$\Gamma \fCenter \Delta, !B$
\AxiomC{}
\RightLabel{$\pi_2'$}
\Deduce$!B, \Gamma \fCenter \Delta$
\RightLabel{\CutCS}
\BinaryInf$\Gamma \fCenter \Delta$
\DisplayProof
\]
has cut rank $= \depth{!B} < \depth{!B \lor !C}$, so the inductive
hypothesis applies again, and yields !!a{proof} $\pi'$ of $\Gamma
\fCenter \Delta$.

\item $!A \ident (!B \land !C)$. Exercise.

\item If $!A \ident (!B \lif !C)$ then $\pi$ ends in
\[
\AxiomC{}
\RightLabel{$\pi_1'$}
\Deduce$!B, \Gamma \fCenter \Delta, !C$
\RightLabel{\RightR{\lif}}
\UnaryInf$\Gamma \fCenter \Delta, !B \lif !C$
\LeftSubproofLabel{$\pi_1$}
\AxiomC{}
\RightLabel{$\pi_2'$}
\Deduce$\Gamma \fCenter \Delta, !B$
\AxiomC{}
\RightLabel{$\pi_2''$}
\Deduce$!C, \Gamma \fCenter \Delta$
\RightLabel{\LeftR{\lif}}
\BinaryInf$!B \lif !C, \Gamma \fCenter \Delta$
\RightSubproofLabel{$\pi_2$}
\RightLabel{\CutCS}
\BinaryInf$\Gamma \fCenter \Delta$
\DisplayProof
\]
The !!{proof} ending in \CutCS,
\[
\AxiomC{}
\RightLabel{$\pi_1'$}
\Deduce$!B, \Gamma \fCenter \Delta, !C$
\AxiomC{}
\RightLabel{$\pi_2''$}
\Deduce$!C, \Gamma \fCenter \Delta$
\RightLabel{\CutCS}
\BinaryInf$!B, \Gamma \fCenter \Delta$
\DisplayProof
\]
has lower cut rank than~$\pi$. The inductive hypthesis yields
!!a{proof}~$\pi_1''$ of $!B, \Gamma \fCenter \Delta$.
Now consider the !!{proof} ending in \CutCS,
\[
\AxiomC{}
\RightLabel{$\pi_2'$}
\Deduce$\Gamma \fCenter \Delta, !B$
\AxiomC{}
\RightLabel{$\pi_1''$}
\Deduce$!B, \Gamma \fCenter \Delta$
\RightLabel{\CutCS}
\BinaryInf$\Gamma \fCenter \Delta$
\DisplayProof
\]
Its cut rank is also less than the cut rank of~$\pi$, so the induction
hypothesis yields a \Log{G3c}-!!{proof}~$\pi'$ of $\Gamma \fCenter
\Delta$.
\item If $!A \ident \lforall[x][!B(x)]$ then $\pi$ ends in
\[
\AxiomC{}
\RightLabel{$\pi_1'$}
\Deduce$\Gamma \fCenter \Delta, !B(c)$
\RightLabel{\RightR{\lforall}}
\UnaryInf$\Gamma \fCenter \Delta, \lforall[x][!B(x)]$
\LeftSubproofLabel{$\pi_1$}
\AxiomC{}
\RightLabel{$\pi_2'$}
\Deduce$!B(t), \lforall[x][!B(x)], \Gamma \fCenter \Delta$
\RightLabel{\LeftR{\lforall}}
\UnaryInf$\lforall[x][!B(x)], \Gamma \fCenter \Delta$
\RightSubproofLabel{$\pi_2$}
\RightLabel{\CutCS}
\BinaryInf$\Gamma \fCenter \Delta$
\DisplayProof
\]
By inductive hypothesis, \CutCS{} can be eliminated from
\[
\AxiomC{}
\RightLabel{$\pi_1''$}
\Deduce$!B(t), \Gamma \fCenter \Delta, \lforall[x][!B(x)]$
\AxiomC{}
\RightLabel{$\pi_2'$}
\Deduce$!B(t), \lforall[x][!B(x)], \Gamma \fCenter \Delta$
\RightLabel{\CutCS}
\BinaryInf$!B(t), \Gamma \fCenter \Delta$
\DisplayProof
\]
where $\pi_1''$ is obtained from $\pi_1$ (not~$\pi_1'$!) by
!!{height}-preserving weakening. (The !!{proof} has the same cut rank
but lower cut !!{height}.) This gives us !!a{proof}~$\pi_2''$ of
$!B(t), \Gamma \fCenter \Delta$.

The !!{proof}~$\pi_1'$ has end-sequent $\Gamma \Sequent \Delta,
!B(c)$, where $c$~is the eigenvariable of a~$\RightR{\lforall}$
inference. So $c$~does not occur in~$!B(x)$, $\Gamma$, or~$\Delta$. We
assume the !!{proof}~$\pi$ is regular, so $c$~is the eigenvariable of
only the following~$\RightR{\lforall}$ inference and only occurs above
it, i.e., only occurs in~$\pi_1'$ and is not an eigenvariable of an
inference in~$\pi_1'$. Since $c$ does not occur in~$\pi_2$, it cannot
occur in~$t$. By \olref[seq][qua]{lem:subst-regular}, the
!!{proof}~$\Subst{\pi_1'}{t}{c}$ is !!a{proof} of $\Gamma \Sequent \Delta,
!B(t)$. We now apply the inductive hypothesis to the !!{proof}
\[
\AxiomC{}
\RightLabel{$\Subst{\pi_1'}{t}{c}$}
\Deduce$\Gamma \fCenter \Delta, !B(t)$
\AxiomC{}
\RightLabel{$\pi_2''$}
\Deduce$!B(t), \Gamma \fCenter \Delta$
\RightLabel{\CutCS}
\BinaryInf$\Gamma \fCenter \Delta$
\DisplayProof
\]
(The induction hypothesis applies since the cut !!{formula}
is~$!B(t)$, so the !!{proof} has lower cut rank than~$\pi$.)
\item We leave the case $!A \ident \lexists[x][!B(x)]$ as an exercise.
\end{enumerate}
\end{proof}

\begin{prob}
Verify the sub-case of case~D in the proof of \olref{lem:cut-adm-G3c}
where $!A$ is not principal in the last inference of~$\pi_1$, which
ends in~$\LeftR{\lif}$. (Note: part of this exercise is to clearly
state what you have to prove, i.e., what form $\pi$ has in this case.)
\end{prob}

\begin{prob}
Verify the sub-case of case~D in the proof of \olref{lem:cut-adm-G3c}
where $!A$ is not principal in the last inference of~$\pi_1$, which
ends in~$\LeftR{\lforall}$.
\end{prob}

\begin{prob}
Verify the sub-case of case~E in the proof of \olref{lem:cut-adm-G3c}
where $!A$ is principal in the last inferences of both $\pi_1$ and
$\pi_2$ and $!A \ident (!B \land !C)$.
\end{prob}

\begin{prob}
Verify the sub-case of case~E in the proof of \olref{lem:cut-adm-G3c}
where $!A$ is principal in the last inferences of both $\pi_1$ and
$\pi_2$ and $!A \ident \lexists[x][!B(x)]$.
\end{prob}

Note that in case~E the cut !!{height} of the !!{proof} ending in
\CutCS{} to which we apply the inductive hypothesis can be larger than
the cut !!{height} of the original proof. E.g., in the case where $!A
\ident (!B \lif !C)$ we turned $\pi$ into !!a{proof} involving two
\CutCS{} inferences,
\[
\AxiomC{}
\RightLabel{$\pi_2'$}
\Deduce$\Gamma \fCenter \Delta, !B$
\AxiomC{}
\RightLabel{$\pi_1'$}
\Deduce$!B, \Gamma \fCenter \Delta, !C$
\AxiomC{}
\RightLabel{$\pi_2''$}
\Deduce$!C, \Gamma \fCenter \Delta$
\RightLabel{\CutCS}
\BinaryInf$!B, \Gamma \fCenter \Delta$
\RightSubproofLabel{$\pi_1''$}
\RightLabel{\CutCS}
\BinaryInf$!B, \Gamma \fCenter \Delta$
\DisplayProof
\]
The top \CutCS{} between $\pi_1'$ and $\pi_2''$ has both lower cut
rank and lower cut !!{height} than~$\pi$. But the
!!{proof}~$\pi_1''$---the result of applying the inductive
hypothesis---no longer has lower cut !!{height} than~$\pi$. Since the
cut !!{formula} of the lower \CutCS{} is $!B$, its cut \emph{rank},
however, is lower than the cut rank of~$\pi$.

\olref{lem:cut-adm-G3c} establishes that we can eliminate a single topmost
\CutCS{} inference from !!a{proof}. We can show that we can eliminate
any number of \CutCS{} inferences from !!a{proof} by successively
applying it to the topmost sub-!!{proof}s ending in~\CutCS.

\begin{thm}\ollabel{thm:cut-elim-G3c-topmost} Any !!{proof} $\pi$ in
$\Log{G3c} + \CutCS$ can be transformed into !!a{proof} in~\Log{G3c}.
\end{thm}

\begin{proof}
  By induction on the number~$n$ of \CutCS{} inferences in~$\pi$. If
  $n=0$, there is nothing to do: $\pi$ is already !!a{proof}
  in~\Log{G3c}. Otherwise, consider the sub-!!{proof}~$\pi'$ ending in
  a top-most \CutCS{} inference. It contains a single \CutCS{} rule as the
  last inference. Apply \olref{lem:cut-adm-G3c} to transform it into a
  \CutCS{}-free proof~$\pi''$, and replace the sub-!!{proof}~$\pi'$
  by~$\pi''$. The result is !!a{proof} containing $n-1$ \CutCS{}
  inferences, and the induction hypothesis applies.
\end{proof}

\end{document}